\documentclass[10pt]{article}

\usepackage[margin=1in]{geometry}
\usepackage{booktabs}
\usepackage[hidelinks]{hyperref}
\usepackage{microtype}
\usepackage{tabularx}

\title{Can an Agent Memory Layer Prove What the Agent Knew?\\
An Evaluation of Regulated-Memory Invariants}
\author{Ethan Beirne\\Lians\\\texttt{https://www.lians.ai}}
\date{July 2026}

\begin{document}
\maketitle

\begin{abstract}
Long-term agent-memory benchmarks primarily measure whether a system can retrieve
useful evidence from conversational or interaction histories. Regulated uses add a
different requirement: an operator may need to reconstruct what an agent knew at a
past decision time, demonstrate that later facts could not influence that decision,
and prove that erased material is no longer recoverable. We introduce an open
evaluation of five regulated-memory invariants: stale-revision suppression,
point-in-time recall, provable erasure, lookahead protection, and historical
audit-state reconstruction. The harness evaluates product-level primitives rather
than general architecture quality. Lians, mem0 OSS, and Graphiti OSS are executed
live; Letta, Hindsight, and Supermemory are capability-assessed from their public API
surfaces pending vendor-supplied live configurations. Under the stated scoring rule,
Lians scores 5.0 of 5, Graphiti 2.0, Letta 1.0, Hindsight 1.0, Supermemory 1.0, and
mem0 OSS 0.5. These results do not establish general memory superiority. They expose
a distinct evaluation axis and provide runnable adapters, per-cell evidence, and a
public correction process for testing it.
\end{abstract}

\section{Introduction}

Persistent memory is becoming a standard component of long-running agents.
LoCoMo and LongMemEval test long-horizon conversational recall
\cite{locomo2024,longmemeval2024}, while LongMemEval-V2 evaluates whether agents can
internalize environment-specific experience \cite{longmemevalv2_2026}. These tasks
are important, but they do not directly test whether a deployed memory layer can
serve as evidence about a past agent decision.

That distinction matters when facts change. A system can answer the current question
correctly while being unable to reconstruct the knowledge state used yesterday. It
can delete a row while being unable to produce an erasure artifact. It can retrieve a
highly relevant fact that was not knowable at the simulated decision time. Recent
work has made state coordination, contradiction resolution, and bitemporal operators
first-class memory problems \cite{toki2026,worlddb2026,atma2026}. We complement that
work with a product-facing evaluation of auditable operational primitives.

This paper contributes:

\begin{enumerate}
  \item Five falsifiable regulated-memory invariants with explicit pass, partial,
        and absent criteria.
  \item A runnable adapter interface that separates live behavioral evidence from
        capability assessment.
  \item An initial six-system comparison with per-cell evidence and public vendor
        right of reply.
  \item A methodology for evaluating what an agent could have known at a past time,
        rather than only whether it can answer correctly now.
\end{enumerate}

\section{Evaluation Contract}

The evaluation treats each invariant as a turnkey product primitive. A pass requires
a supported API path that satisfies the full invariant. Partial credit indicates a
related capability that leaves required filtering, proof construction, or state
reconstruction to the application. Absent means that no supported primitive was
identified. Scores are one point for pass, one half point for partial, and zero for
absent.

\subsection{Invariants}

\paragraph{Stale-revision suppression.}
After a keyed fact is revised, default retrieval must not return the superseded value
as current context. Preserving old values for history is compatible with a pass if
the retrieval contract distinguishes them deterministically.

\paragraph{Point-in-time recall.}
The caller can query the knowledge state as known at a prior transaction time. This
is distinct from filtering only on when a fact was valid in the external world.

\paragraph{Provable erasure.}
Erased content becomes cryptographically unrecoverable and the system emits a proof
artifact or stable request reference. Behavioral non-retrieval without a proof
artifact receives partial credit.

\paragraph{Lookahead protection.}
The system prevents or flags retrieval of facts that were not knowable at a declared
simulation or decision time.

\paragraph{Historical audit-state reconstruction.}
The caller can reproduce the complete knowledge state that was available at a past
time, including the effect of revisions and erasures under the documented policy.

\section{Systems and Execution Status}

Table~\ref{tab:status} distinguishes behavioral runs from capability assessment. This
distinction is central. A capability-assessed cell is a documented hypothesis about
the public product surface, not a claim that the system was executed.

\begin{table}[h]
\centering
\caption{Evidence status by system.}
\label{tab:status}
\begin{tabular}{lll}
\toprule
System & Evidence mode & Configuration \\
\midrule
Lians & Live & LocalLiansClient, version 0.4.0 \\
mem0 & Live & OSS 2.0.11, documented OpenAI configuration \\
Graphiti & Live & OSS 0.29.2, OpenAI clients, embedded Kuzu \\
Letta & Capability-assessed & Public API and adapter \\
Hindsight & Capability-assessed & Public API and adapter \\
Supermemory & Capability-assessed & Public API and adapter \\
\bottomrule
\end{tabular}
\end{table}

The live OSS configurations follow vendor quickstarts. Any accommodation required to
complete a run is recorded in the public results appendix. Hosted adapters switch to
live execution when the corresponding credentials are supplied. The harness does not
convert an unexecuted capability map into behavioral evidence.

\section{Initial Results}

\begin{table*}[t]
\centering
\caption{Initial regulated-memory results. P indicates pass, H indicates half credit,
and A indicates absent. Capability-assessed columns remain subject to live vendor
configuration.}
\label{tab:results}
\begin{tabularx}{\textwidth}{Xcccccc}
\toprule
Invariant & Lians & Graphiti & Letta & Hindsight & Supermemory & mem0 \\
\midrule
Stale revision suppressed & P & H & H & H & H & A \\
Point-in-time recall & P & H & A & H & A & A \\
Provable erasure & P & H & H & A & H & H \\
Lookahead protection & P & A & A & A & A & A \\
Audit-state snapshot & P & H & A & A & A & A \\
\midrule
Score & 5.0 & 2.0 & 1.0 & 1.0 & 1.0 & 0.5 \\
\bottomrule
\end{tabularx}
\end{table*}

The results identify a product gap rather than a universal ranking. Graphiti's
contradiction invalidation correctly backdated stale edges in the live run, but
default retrieval returned invalidated edges unless the caller filtered them. mem0
made erased content behaviorally unretrievable but emitted no proof artifact. These
capabilities receive partial credit under the same rule applied to Lians and all
other systems.

\section{Reproducibility and Correction Process}

The implementation, adapters, generated table, and per-cell notes are available in
the Lians repository \cite{lians2026}. A basic run is:

\begin{verbatim}
cd agentmem
python -m benchmarks.compare_regulated
\end{verbatim}

Optional vendor credentials activate hosted adapters. The public results document
records exact OSS versions and accommodations. Each evaluated vendor received a
public request to identify missed APIs, correct characterizations, or provide a
supported configuration. Accepted corrections change the adapter and the published
result, whether they improve or reduce a vendor's score.

\section{Threats to Validity}

\paragraph{Evaluator conflict of interest.}
The evaluation is authored by the developer of Lians. Open adapters, public evidence,
and vendor correction requests reduce but do not eliminate this conflict. Independent
replication is required.

\paragraph{Mixed evidence modes.}
Three columns are capability-assessed rather than live. They must not be interpreted
as behavioral measurements. The primary near-term objective is to replace these
columns with supported live configurations.

\paragraph{Configuration sensitivity.}
Memory behavior can change with model, prompt, version, storage backend, and hosted
product tier. Results apply only to the recorded configuration.

\paragraph{Binary product criteria.}
Turnkey primitives favor systems that expose explicit APIs over architectures from
which an application could construct equivalent behavior. This is intentional for
procurement analysis, but it is not a measure of theoretical expressiveness.

\paragraph{No utility or latency claim.}
The evaluation does not measure personalization quality, answer accuracy, cost, or
latency. A system that scores lower here may be superior for many agent workloads.

\section{Related Work}

LoCoMo and LongMemEval established long-horizon conversational evaluation
\cite{locomo2024,longmemeval2024}. LongMemEval-V2 shifts attention toward memory of
environment experience \cite{longmemevalv2_2026}. Look-Ahead-Bench evaluates
model-level temporal leakage in financial workflows by comparing performance and
alpha decay across market regimes \cite{lookaheadbench2026}. Our lookahead invariant
is complementary: it tests whether a memory product prevents facts that were not
knowable at the declared decision time from entering retrieval context, even when
the underlying model is temporally clean. WorldDB evaluates immutable,
content-addressed worlds and write-time reconciliation \cite{worlddb2026}. TOKI
formalizes contradiction resolution as bitemporal concurrency control with provenance
contracts \cite{toki2026}. A-TMA separates bank maintenance, retrieval, and answer-time
state resolution, exposing ghost-memory failures hidden by final answer accuracy
\cite{atma2026}. Our evaluation operates at a complementary layer: whether a product
surface exposes auditable primitives that a regulated operator can execute directly.

\section{Conclusion}

Agent memory has two different evaluation questions. One asks whether the agent can
recall useful information. The other asks whether an operator can prove which
information was available, prevent future information from leaking backward, and
produce evidence about revision and erasure. The initial results suggest that the
second question is not captured by current recall benchmarks or most default product
surfaces. The proposed invariants, open harness, and correction process provide a
starting point for independent evaluation rather than a final ranking.

\bibliographystyle{plain}
\bibliography{references}

\end{document}
